%%%PAGE 211 rot=0 legibility=good summary=测电源E、r实验：方程组法(伏安/伏阻/安阻)、补偿电路消系统误差、定值电阻的应用、内外接改进与修正公式、特殊伏阻法与特殊安阻法
%%%BLOCK id=211-01 ch=10 sec=测电动势和内阻·方程组法 kind=formula exp=1 alt=none cont=none y=0.03-0.19
\begin{center}
\renewcommand{\arraystretch}{1.6}
\begin{tabular}{@{}c c c c@{}}
 & 伏安法 & 伏阻法 & 安阻法\\
方程组法 &
$\left\{\begin{array}{l} E=\underline{U_1}+\underline{I_1}r\\ E=\underline{U_2}+\underline{I_2}r\end{array}\right.$ &
$\begin{array}{l} U_1=E\cdot\dfrac{R_1}{R_1+r}\\[6pt] U_2=E\,\dfrac{R_2}{R_2+r}\end{array}$ &
$\begin{array}{l} E=I_1(R_1+r)\\ E=I_2(R_2+r)\end{array}$
\end{tabular}
\end{center}
%%%END
%%%BLOCK id=211-02 ch=10 sec=测电动势和内阻·消除系统误差 kind=method exp=1 alt=none cont=none y=0.23-0.37
%FIG p211_a 0.08 0.225 0.39 0.365
\notefig[0.31]{figures/p211_a.png}

$i_1$ $i_2$ 抵消。无误差（\unclear{除了表所带系统误差}）

$E_1=E_2$

$r_1=r_2$
%%%END
%%%BLOCK id=211-03 ch=10 sec=测电动势和内阻·定值电阻的应用 kind=method exp=1 alt=none cont=none y=0.39-0.57
定值电阻在测电源 $E$、$r$ 实验中的应用。

%FIG p211_b 0.07 0.455 0.60 0.565
\notefig[0.5]{figures/p211_b.png}

图旁标注（自上而下，依次在 $R_1$、$R_2$、$R_3$ 所在行右侧）：
\begin{itemize}
\item 看成滑变的一部分。
\item 改装电表
\item 看成电源内阻的一部分
\end{itemize}
%%%END
%%%BLOCK id=211-04 ch=10 sec=测电动势和内阻·内外接改进 kind=method exp=1 alt=none cont=none y=0.57-0.71
\textbf{改进}（左侧大括号）：
\begin{itemize}
\item 内外均\unclear{结}法，
\item \underline{内接$\to E$准确} % 红笔划线
\item \underline{外接$\to I_0$准确} \quad \underline{$I_0=\dfrac{E}{r}\to r$ 准确值} % 红笔划线（两段：“外接…准确”与“$I_0=E/r\to r$准确值”）
\item \underline{A内V外法} % 红笔波浪线
\end{itemize}
%%%END
%%%BLOCK id=211-05 ch=10 sec=测电动势和内阻·内外接修正公式 kind=formula exp=1 alt=none cont=none y=0.72-0.79
\[
\text{已知 }R_A\to\text{内接}\
\left\{\begin{array}{l} E_{\text{测}}=E_{\text{真}}\\ r_{\text{测}}=r_{\text{真}}+R_A \end{array}\right.
\qquad
\text{已知 }R_V\to\text{外接}\
\left\{\begin{array}{l} E_{\text{测}}=E_{\text{真}}\dfrac{R_V}{R_V+r}\\[8pt] r_{\text{测}}=\dfrac{r_{\text{真}}R_V}{R_V+r_{\text{真}}} \end{array}\right.
\]
%%%END
%%%BLOCK id=211-06 ch=10 sec=测电动势和内阻·特殊伏阻法 kind=method exp=1 alt=none cont=none y=0.83-1.00
\textbf{一种特殊的伏阻法}

%FIG p211_c 0.04 0.862 0.545 0.955
\notefig[0.5]{figures/p211_c.png}

（箭头上方：\unclear{看作…}；箭头下方：\unclear{就在电源里}）

\[
U=\frac{E R_0}{R_0+r+R}
\qquad
\frac{1}{U}=\frac{1}{ER_0}\cdot R+\frac{R_0+r}{ER_0}
\]
%%%END
%%%BLOCK id=211-07 ch=10 sec=测电动势和内阻·特殊安阻法 kind=method exp=1 alt=none cont=none y=0.83-1.00
\textbf{一种特殊的安阻法}

%FIG p211_d 0.60 0.865 0.762 0.932
\notefig[0.25]{figures/p211_d.png}

%FIG p211_e 0.765 0.865 0.925 0.958
\notefig[0.25]{figures/p211_e.png}

\[ E=I(R+r) \]
\[ IR=-r\cdot I+E \]
%%%END
%%%PAGEEND 211
